\documentclass[12pt,a4paper]{report}

\usepackage[utf8]{inputenc}
\usepackage[T1]{fontenc}
\usepackage[english]{babel}
\usepackage{array}
\usepackage{float}
\usepackage{geometry}
\usepackage{longtable}
\usepackage{listings}
\usepackage{fvextra}
\usepackage{csquotes}
\usepackage[hyphens]{url}
\usepackage[hidelinks,breaklinks=true]{hyperref}

\geometry{margin=2.5cm}
\setlength{\parindent}{0pt}
\setlength{\parskip}{0.6em}
\renewcommand{\arraystretch}{1.2}

\lstset{
  basicstyle=\small\ttfamily,
  breaklines=true,
  breakatwhitespace=false,
  columns=fullflexible,
  frame=single
}

\title{LifeRouter\\Technical Manual}
\author{Kentalives}
\date{June 2026}

\begin{document}

\pagenumbering{roman}
\maketitle
\setcounter{page}{2}
\tableofcontents
\clearpage
\pagenumbering{arabic}

\chapter{Technical Manual}

This manual contains the information needed to modify, test, and maintain LifeRouter.
 The user manual describes installation, configuration,
and use of the public API. This document focuses on the internal code structure,
package dependencies, service lifecycle, and the points where developers must be
careful before introducing changes.

\section{Internal System Organization}

The project is a Go module that combines a public API for integrations, an
executable entry point for running the service as a standalone process, and an
\texttt{internal/} tree reserved for implementation code. This separation matters
because external consumers can import public packages, while internal packages
must only be used by this module.

\begin{table}[H]
\centering
\begin{tabular}{|p{0.31\textwidth}|p{0.60\textwidth}|}
\hline
\textbf{Path} & \textbf{Main responsibility} \\
\hline
\texttt{main.go} &
Standalone process entry point. It loads configuration, creates the default
dependencies, runs the dispatcher, and waits for system shutdown signals. \\
\hline
\texttt{embedded/} &
Public surface for running the service inside another Go process. It exposes
default configuration, default dependency wiring, and the startup call that
delegates to the same application path as the binary. \\
\hline
\texttt{pkg/config} &
Public configuration and dependency-injection types. Changes to these types can
break external repositories that build custom configurations. \\
\hline
\texttt{pkg/domain} &
Exported domain types: costs, visualization states, emergency directions,
registered NATS-domain errors, and geospatial query or publishing interfaces. \\
\hline
\texttt{pkg/pathfinding} and \texttt{pkg/emergency} &
Go clients over NATS for starting agent routes, controlling communicators, and
managing emergency routing. They do not own the NATS connection they receive. \\
\hline
\texttt{pkg/subjects} &
Public constants for the NATS subjects handled by the service. They allow direct
NATS use without duplicating string literals in Go consumers. \\
\hline
\texttt{pkg/}\newline\texttt{snapshotstream} &
Shared codec for per-floor snapshots sent through NATS streams. It preserves the
simple shape for small messages and chunks large snapshots with
\texttt{gzip+json}. \\
\hline
\texttt{pkg/tuning} &
Public helper tools for inspecting rasterized maps and tuning the grid before
running the service, including textual output and a browser visualizer for large
grids. \\
\hline
\end{tabular}
\caption{Executable entry point and public API responsibilities.}
\end{table}

\begin{table}[H]
\centering
\begin{tabular}{|p{0.31\textwidth}|p{0.60\textwidth}|}
\hline
\textbf{Path} & \textbf{Main responsibility} \\
\hline
\texttt{internal/app} &
Application wiring, world initialization, NATS dispatcher, request handlers,
active-agent state, and emergency lifecycle. \\
\hline
\texttt{internal/config} &
Configuration loading with Viper, default values, and process-wide storage for
configuration and dependencies used by internal packages. \\
\hline
\texttt{internal/grid} and \texttt{internal/raster} &
Grid representation, coordinates, multi-floor world, portals, and conversion
from map images into traversal costs. \\
\hline
\texttt{internal/pathfinding} &
Algorithm implementations: shared primitives, dynamic agent planning, and
emergency flow fields. \\
\hline
\texttt{internal/mock}, \texttt{internal/} \newline \texttt{pathfinding/testing}, and
\texttt{internal/integration} &
Test doubles, shared scenarios, and global behavior tests for the service. \\
\hline
\texttt{data/} &
Map images used as inputs in tests and rasterization tuning. \\
\hline
\end{tabular}
\caption{Internal implementation packages and supporting data.}
\end{table}

\subsection{Boundary Between Public API and Implementation}

The public boundary is formed by \texttt{embedded/} and \texttt{pkg/}. The
\texttt{embedded} package lets another Go application run the service in its own
process through \texttt{DefaultConfig}, \newline \texttt{DefaultDependencies}, and
\texttt{Run}. This path is useful for integrations, tools, and debugging runs,
but it still requires a NATS connection and domain dependencies that are
coherent with the real environment.

The \texttt{pkg/} directory contains the types and clients that other
repositories can use. The \texttt{pkg/pathfinding} and \texttt{pkg/emergency}
packages translate Go calls into NATS requests. \texttt{pkg/config} defines
configuration and injectable dependencies. \texttt{pkg/domain} contains shared
types such as costs, directions, preferred-route graphs, and domain errors.
\newline \texttt{pkg/subjects} exposes NATS subject constants for consumers that build
messages directly. When this surface changes, both internal tests and possible
external integrations must be considered.

By contrast, \texttt{internal/} must not be treated as a stable API. It contains
global application state, algorithms, and communication details. A change inside
\texttt{internal/} may legitimately avoid external compatibility, but it must
preserve the behavior observed through the public API and NATS subjects.

\subsection{Service Lifecycle}

The standalone process and the embedded mode use the same application path. The
main difference is who creates the configuration, dependencies, and cancellation
context.

\begin{table}[H]
\centering
\begin{tabular}{|p{0.24\textwidth}|p{0.67\textwidth}|}
\hline
\textbf{Phase} & \textbf{Technical description} \\
\hline
Initial load &
\texttt{main.go} creates a context cancellable by \texttt{SIGINT} or
\texttt{SIGTERM}, calls \texttt{config.LoadConfig}, and obtains values from
defaults, \texttt{config/config.yaml}, and environment variables. \\
\hline
Default dependencies &
\texttt{embedded.DefaultDependencies} opens separate NATS connection pools for
geospatial queries and publishing, and prepares a simulated external system.
Production deployments can provide their own dependencies through
\texttt{pkg/config.Dependencies}. \\
\hline
Application startup &
\texttt{app.Run} stores global configuration and dependencies, initializes
logging, builds the rasterized world with layers and portals, creates the
service NATS connection, and registers handlers. \\
\hline
Normal execution &
The \texttt{Dispatcher} keeps active agent communicators, briefly retains
completed agent results, and manages the current emergency state. Blocking
request paths use worker limits to avoid unbounded goroutine growth. \\
\hline
Graceful shutdown &
\texttt{Dispatcher.Shutdown} stops active agents, stops the emergency if one is
running, drains the service NATS connection, cancels the service context, and
closes dependency resources that expose \texttt{Close}. \\
\hline
\end{tabular}
\caption{Main lifecycle phases of the microservice.}
\end{table}

\subsection{Dependency Model}

Internal packages access configuration and dependencies through
\texttt{internal/config}. This simplifies access from algorithms and handlers,
but it requires service startup to call \texttt{config.Setup} before code that
reads \texttt{config.Cfg} or \texttt{config.Dep}. Tests that exercise internal
packages must explicitly prepare this state when they call functions that depend
on it.

\begin{table}[H]
\centering
\begin{tabular}{|p{0.34\textwidth}|p{0.58\textwidth}|}
\hline
\textbf{Element} & \textbf{Use inside the system} \\
\hline
\texttt{pkg/config.Config} &
Complete configuration: NATS connection, logging, grid dimensions, layers,
portals, agent parameters, emergency parameters, and geotruth pool sizes. \\
\hline
\texttt{pkg/config.Dependencies} &
Groups the external systems required for domain data, geospatial reads, and
movement publishing. \\
\hline
\texttt{domain.ExternalSystem} &
Provides information not owned by pathfinding: object costs, floor heights, and
emergency signaling nodes. \\
\hline
\texttt{domain.GeoQuery} &
Read-side geospatial interface: nearby objects, object data, oriented object
sets, and floor/region lookup for a real-world point. \\
\hline
\texttt{domain.GeoPublish} &
Write-side interface used to publish an agent's new position and receive the
external commit acknowledgement. \\
\hline
\end{tabular}
\caption{Configuration and external dependencies used by the application.}
\end{table}

\subsection{Geotruth Connection Pools}

The default dependencies wrap \texttt{GeoQuery} and \texttt{GeoPublish} in
internal NATS client pools. The query pool is used for object data, nearby
objects, oriented objects, and floor resolution. The publish pool is used for
position updates. Separating both flows avoids making local agent vision and
real-position publishing always compete for the same NATS connection.

Pool sizes are configured under \texttt{pathfinding.geotruth} through
\texttt{queryconnections} and \texttt{publishconnections}. If omitted, both
default to \texttt{GOMAXPROCS}; values lower than one are normalized to one
connection. Client selection uses round-robin atomic counters. During shutdown,
\texttt{Dispatcher.Shutdown} closes dependencies that implement \texttt{Close},
including the pools created by \texttt{DefaultDependencies}.

\subsection{NATS Handler Ownership}

The concrete NATS subjects are documented in the user manual and defined as
public constants in \texttt{pkg/subjects}. For maintenance, the important point
is which dispatcher area owns each handler group and which internal state it
changes.

\begin{table}[H]
\centering
\begin{tabular}{|p{0.28\textwidth}|p{0.63\textwidth}|}
\hline
\textbf{Group} & \textbf{Internal responsibility} \\
\hline
Agent pathfinding &
Starts planning for an existing object, converts the goal from meters to grid
coordinates, creates an internal \texttt{AgentCommunicator}, and replaces any
previous run for the same agent. \\
\hline
Agent path observation &
Associates a state-publishing channel with the active communicator and streams
snapshots until the route ends or the communicator closes the channel. \\
\hline
Agent communicator control &
Checks movement state, waits for completion, changes speed, requests discrete or
meter-based movement, pauses automatic movement, or terminates the route. \\
\hline
Emergency pathfinding &
Applies preferred routes to the global world, converts goals to grid
coordinates, starts the flow field, and prevents simultaneous emergencies. \\
\hline
Emergency flow observation &
Publishes per-floor direction maps while an emergency is active and closes the
subscription when the flow publisher finishes. \\
\hline
Emergency stop and shutdown &
Coordinates emergency shutdown, waits for completion, and restores global costs
when preferred routes are removed. \\
\hline
\end{tabular}
\caption{NATS handler groups and associated internal state.}
\end{table}

\subsection{Validation Points for Changes}

Changes to public APIs or NATS messages must be reviewed together with the user
manual because they alter the integration contract. Changes to the dispatcher or
lifecycle should run the \texttt{internal/app} tests and integration tests.
Changes to grid logic, rasterization, or algorithms require the corresponding
unit tests and, if performance is affected, the profiling or benchmark tests
already separated in the repository.

The general validation command is:

\begin{lstlisting}
make test
\end{lstlisting}

During local iteration, run the package-level tests first. Before closing a
change, run the full suite. When measuring code exercised by all tests across
packages, use the full coverage target described in the testing section.

\section{Spatial Representation: Grid, Rasterization, and Multi-Floor World}

The spatial representation is shared by agent pathfinding and emergency routing.
The service does not work directly on continuous geometry. It works on a
per-floor cost grid. Public API coordinates arrive in meters; before algorithms
run, those coordinates are converted to cells and associated with a world layer.

\subsection{Units and Coordinates}

Internal local coordinates use \texttt{grid.Coords}, which contains integer
\texttt{X} and \texttt{Y} cell indexes. Conversion between meters and cells
depends on \texttt{GridConfig.CellSizeM}. Cell-to-meter conversion returns the
cell center. Meter-to-cell conversion rounds toward the corresponding cell and
never returns negative coordinates.

\begin{table}[H]
\centering
\begin{tabular}{|p{0.34\textwidth}|p{0.58\textwidth}|}
\hline
\textbf{Element} & \textbf{Behavior to preserve} \\
\hline
\texttt{CellSizeM} &
Real-world side length represented by one cell. It must be greater than zero;
invalid global configuration causes the grid package to panic. \\
\hline
\texttt{Coords.ToFloat64} &
Converts cell indexes to meters and returns the center of the cell. It is the
reference for publishing or comparing real positions derived from the grid. \\
\hline
\texttt{CoordsFromFloat64} &
Converts an \texttt{x,y} meter position to a cell coordinate using active
configuration and clamps negative coordinates to zero. \\
\hline
\texttt{GlobalCoords} &
Adds a layer index to a local coordinate, distinguishing the same \texttt{x,y}
cell on different floors. \\
\hline
\texttt{GlobalCoordsFromFloat64} &
Uses \texttt{GeoQuery.RegionFromPoint} to resolve the layer and then converts
\texttt{x,y} to a cell. It requires \texttt{config.Dep.Qu} to be initialized. \\
\hline
\end{tabular}
\caption{Main units and conversions in the spatial representation.}
\end{table}

Global coordinates must be validated before converting them to flattened
indexes. A cell whose \texttt{X} or \texttt{Y} is outside its floor cannot rely
on \texttt{ToIdx}, because arithmetic could alias another floor's index range.
The intended exception is the emergency fake goal, represented as a sentinel
outside the grid with flattened index \texttt{-1}.

\subsection{Grid Cost Layers}

Each \texttt{Grid} keeps three same-size buffers. The base layer comes from the
rasterized map and floor masks. The object layer represents dynamic costs, such
as obstacles or observed agents. The final layer is the effective cost read by
the algorithms. This separation keeps the base map stable while updating only
regions affected by dynamic information.

\begin{table}[H]
\centering
\begin{tabular}{|p{0.31\textwidth}|p{0.61\textwidth}|}
\hline
\textbf{Concept} & \textbf{Behavior to preserve} \\
\hline
\texttt{base} &
Permanent floor costs. A zero value blocks a cell and must not be overwritten by
ordinary dynamic updates. \\
\hline
\texttt{objects} &
Temporary costs added by objects or local updates. A zero-cost object region can
clear previous temporary costs and restore the base cost. \\
\hline
\texttt{final} &
Cost used by pathfinding. It is recomputed with \texttt{AddCost(base, objects)}
and change notifications are emitted when it changes. \\
\hline
\texttt{RegionGrid} &
Rectangular region relative to a parent grid. It is used to replace object
costs, apply base masks, and limit updates around a position. \\
\hline
\texttt{SubChanges} &
Registers synchronous callbacks for batches of changed cells. Updates must
preserve notification of previous costs. \\
\hline
\texttt{AddCost} and \texttt{DiagonalCost} &
Cost operations that are safe around unreachable values and overflow. Algorithms
must use these helpers rather than adding costs directly. \\
\hline
\end{tabular}
\caption{Internal grid layers and related operations.}
\end{table}

Main mutation points are object-region replacement, base-mask application,
final-cost recomputation, and base reset. Any change here must preserve blocked
cells, publish each batch of changes once, and keep base, object, and final
layers synchronized. A future extension could page the local object and final
layers, similar to the sparse storage already used for D* Lite values.

\subsection{Map Rasterization}

World creation calls \texttt{grid.FromImg}, which builds each floor base layer
through \newline \texttt{raster.RasterPNG4WithAura}. \texttt{pxpercell} tells how many
source image pixels correspond to one grid cell. Traversable cells are selected
from the red channel: in simple paletted images, the red palette entry is used;
in generic images, the average red value must pass the internal threshold. Wall
aura then increases cost around blocked cells using a Chebyshev neighborhood.

\begin{table}[H]
\centering
\begin{tabular}{|p{0.30\textwidth}|p{0.62\textwidth}|}
\hline
\textbf{Change point} & \textbf{Main risk} \\
\hline
Image format handling &
Changing palette detection or fallback behavior can make existing maps
non-traversable or unexpectedly expensive. \\
\hline
Pixel-to-cell sampling &
Changing center sampling or scaling rules can move walls and corridors relative
to the configured grid. \\
\hline
Wall aura &
Changing the Chebyshev radius or cost can close narrow corridors or remove
clearance around walls. \\
\hline
Base-cost convention &
Changing empty-space and unreachable values affects all grid and pathfinding
cost calculations. \\
\hline
\end{tabular}
\caption{Rasterization risks that require explicit validation.}
\end{table}

\subsection{Multi-Floor World, Portals, and Virtual Worlds}

\texttt{World} groups floor grids and stores metadata that flattens global
coordinates into unique indexes. Layer names are translated to internal indexes
through the configured \newline \texttt{GridConfig.FloorLayers}. Geotruth region names
must match these configured names when public \texttt{[x,y,z]} coordinates are
converted into layers.

Portals are directed cross-floor edges. Global portals are configured during
startup. Local portals belong only to a leased virtual world and are useful for
agent-specific changes. The world stores portals indexed by origin and by
destination, so \texttt{PortalsFrom} and \texttt{PortalsTo} do not need to scan
unrelated portals.

\begin{table}[H]
\centering
\begin{tabular}{|p{0.32\textwidth}|p{0.60\textwidth}|}
\hline
\textbf{Element} & \textbf{Role} \\
\hline
\texttt{NewWorld} &
Rasterizes configured layers, installs the process-wide global world, and later
receives configured portals. \\
\hline
\texttt{NewVirtualWorld} &
Leases a per-agent world that lazily links floors from the global world and keeps
local portals separate. \\
\hline
\texttt{AddPortal} and \texttt{AddBidirectionalPortal} &
Add directed or paired global cross-floor edges. Invalid costs, same-floor
portals, and invalid endpoints are rejected or logged. \\
\hline
\texttt{AddLocalPortal} &
Adds a portal visible only in the leased virtual world, preserving the global
world for other agents and for emergency routing. \\
\hline
\texttt{PortalsFrom} and \texttt{PortalsTo} &
Merge global and local portals from the corresponding origin or destination
indexes. \\
\hline
\end{tabular}
\caption{Multi-floor world elements and portal responsibilities.}
\end{table}

\subsection{Validation of Spatial Changes}

Spatial changes must be validated at the layer they affect: coordinate
conversion, grid costs, rasterization, portals, or virtual-world reuse.

\begin{table}[H]
\centering
\begin{tabular}{|p{0.31\textwidth}|p{0.61\textwidth}|}
\hline
\textbf{Change} & \textbf{Minimum validation} \\
\hline
Coordinate conversion &
Tests for meter-to-cell and cell-to-meter conversion, negative coordinates,
configured cell size, and layer resolution. \\
\hline
Cost-layer mutation &
Tests that blocked base cells remain blocked, final costs update correctly, and
subscribers receive changed cells. \\
\hline
Rasterization &
Tests with representative map fixtures, aura radius, complex-image fallback, and
red-channel threshold behavior. \\
\hline
Portals &
Unit tests for directed and bidirectional portals, same-floor rejection,
incoming-portal lookup, removal, and virtual-world local portals. \\
\hline
Virtual-world reuse &
Tests that reused worlds do not leak local portals, subscriptions, or dynamic
objects from a previous agent run. \\
\hline
\end{tabular}
\caption{Recommended validation for spatial changes.}
\end{table}

\section{Agent Pathfinding With D* Lite}

Agent routing is based on D* Lite \cite{koenig-likhachev-dstar-lite}. The
planner controls movement, local geotruth vision, dynamic cost updates, and
route publication. Each route runs in a goroutine and owns a virtual world
derived from the global world.

\subsection{Lifecycle of an Agent Route}

\begin{table}[H]
\centering
\begin{tabular}{|p{0.26\textwidth}|p{0.65\textwidth}|}
\hline
\textbf{Phase} & \textbf{Description} \\
\hline
Request reception &
The dispatcher decodes \texttt{AgentFindPathParams}, checks shutdown state,
converts the goal to \texttt{GlobalCoords}, and starts internal agent
pathfinding. \\
\hline
Initial state &
The agent position is read from geotruth, converted to grid coordinates, and
associated with a leased virtual world. \\
\hline
Planning &
D* Lite initializes values, priority queue state, and the current goal; it then
computes the shortest path from the current position. \\
\hline
Movement loop &
While the start is not the goal, the communicator decides whether a step can be
taken based on speed, manual commands, and termination signals. \\
\hline
Local update &
Nearby objects are queried from geotruth, rasterized into local grid overlays,
and changed cells trigger D* Lite repairs. \\
\hline
Real publication &
According to \texttt{CellsForRealUpdate}, internal position is converted to
meters and published with \texttt{GeoPublish.UpdateObjectPosition}; floor height
is read from \texttt{ExternalSystem.HeightAtFloorPoint}. \\
\hline
Completion &
The communicator records the terminal error, closes publishers and waiters,
releases the virtual world, and the dispatcher retains the result briefly. \\
\hline
\end{tabular}
\caption{Agent route lifecycle.}
\end{table}

\subsection{D* Lite State}

\begin{table}[H]
\centering
\begin{tabular}{|p{0.28\textwidth}|p{0.64\textwidth}|}
\hline
\textbf{State} & \textbf{Role} \\
\hline
\texttt{g} and \texttt{rhs} &
The D* Lite value pair. Dense mode uses one contiguous store; sparse mode
allocates value pages only where needed. \\
\hline
Priority queue &
Queue with lookup, update, and removal by index. Keys encode D* Lite priority
and depend on the current start. \\
\hline
\texttt{start} and \texttt{goal} &
The current position is stored as coordinates; the goal is stored as a flattened
global index. \\
\hline
World subscription &
Receives changed cells from the virtual world and transforms them into D* Lite
vertex updates. \\
\hline
Portal state &
Successor and predecessor expansion must include same-floor neighbors and
cross-floor portals. \\
\hline
\end{tabular}
\caption{Main state used by the agent D* Lite planner.}
\end{table}

The sentinel index \texttt{-1} is reserved for the emergency fake goal and must
not be treated as a normal agent pathfinding cell.

\subsection{Local Vision and Dynamic Changes}

Agent replanning is driven by local observations. The planner asks geotruth for
nearby objects inside the configured vision radius converted to meters. Object
geometry is rasterized into the agent's virtual-world object layers. Changed
cells are then propagated to D* Lite, which repairs only the affected parts of
the search state. This conversion follows a scanline-style rasterization
strategy commonly used to fill geometry over pixel or cell grids
\cite{foley1995computergraphics,pineda1988polygon}.

\begin{table}[H]
\centering
\begin{tabular}{|p{0.31\textwidth}|p{0.61\textwidth}|}
\hline
\textbf{Element} & \textbf{Meaning} \\
\hline
\texttt{VisionRadiusCells} &
Radius of local geotruth refresh in grid cells; it is also converted to meters
for the external query. \\
\hline
\texttt{SparseValuePageDirectory} &
Enables sparse storage for D* Lite \texttt{g/rhs} values in large local worlds. \\
\hline
Object rasterization &
Converts oriented objects into dynamic grid costs using object traversal costs
from the external system. \\
\hline
Changed-cell repair &
Updates predecessor and successor costs so D* Lite can reuse previous search
work instead of planning from scratch. \\
\hline
\end{tabular}
\caption{Configuration and dependencies used by agent local vision.}
\end{table}

\subsection{Communicator and Movement Control}

The \texttt{AgentCommunicator} connects public commands, internal route state,
and goroutine lifecycle. It must maintain the ordering of commands for one agent;
worker limits in NATS handlers do not parallelize one agent's algorithm.

\begin{table}[H]
\centering
\begin{tabular}{|p{0.30\textwidth}|p{0.62\textwidth}|}
\hline
\textbf{Responsibility} & \textbf{Behavior} \\
\hline
\texttt{UpdateMovementSpeed} &
Changes continuous movement speed in cells per second. A zero value pauses
automatic movement without terminating the job. \\
\hline
\texttt{MoveNCells} &
Requests a bounded number of grid-cell steps. \\
\hline
\texttt{MoveFMeters} &
Blocks until the requested distance in meters is consumed or the agent exits; it
returns remaining meters when the request cannot be fully consumed. \\
\hline
\texttt{Stop} &
Pauses automatic movement by setting speed to zero. \\
\hline
\texttt{Terminate} &
Cancels the route and stores \texttt{ErrAgentTerminated} if it was still active. \\
\hline
\texttt{BlockingWait} and \texttt{ExitError} &
Allow clients to wait for completion and query the terminal domain error. The
dispatcher also supports these calls briefly after completion. \\
\hline
Path publisher &
Streams route-visualization snapshots. Replacing a publisher closes the previous
one. \\
\hline
\end{tabular}
\caption{Agent communicator responsibilities.}
\end{table}

Automatic movement uses absolute next-tick scheduling. The communicator schedules
the next movement instant before returning to route processing, so planning and
publication time count against cadence. If a tick is already overdue, the
implementation permits one immediate step and then resets timing from the current
time rather than catching up with bursts.

\subsection{Resource Ownership and Validation}

An agent run owns its virtual world, local subscriptions, path publisher, wait
channels, and terminal result. Changing this ownership can leave callbacks
active, stale data in reused worlds, or terminal errors unavailable to clients.
Use tests around replacement by a same-agent run, termination, retained results,
and shutdown races whenever this area changes.

\section{Emergency Routing With Flow Fields}

Emergency routing computes a global flow field toward one or more exits. The
flow field is updated when the environment changes and can be mapped to
emergency signaling nodes. Unlike agent routing, this uses the global world
because preferred routes intentionally modify global base costs and are restored
when the emergency stops.

\subsection{Emergency Lifecycle}

\begin{table}[H]
\centering
\begin{tabular}{|p{0.27\textwidth}|p{0.64\textwidth}|}
\hline
\textbf{Phase} & \textbf{Behavior} \\
\hline
Start request &
The dispatcher checks that no emergency is already running, applies preferred
routes, converts goals from meters to grid coordinates, and starts the flow
field. \\
\hline
Initialization &
The emergency planner validates goals, prepares light-node state, adds the fake
goal sentinel, and subscribes to world changes. \\
\hline
Initial computation &
The planner computes the global flow directions for all reachable cells. \\
\hline
Periodic update &
At the configured tick rate, environment objects are refreshed, changed cells
are applied, the D* Lite state is repaired, and flow snapshots are published
when requested. \\
\hline
Signaling update &
Light nodes are updated when cost changes exceed the lighting hysteresis and a
dominant direction is selected from local neighborhood directions. \\
\hline
Stop &
The dispatcher closes the emergency quit channel, waits for completion, removes
preferred-route masks, and marks the emergency as stopped. \\
\hline
\end{tabular}
\caption{Emergency routing lifecycle.}
\end{table}

\subsection{Flow-Field State}

The emergency planner reuses D* Lite ideas, but its output is a direction per
cell rather than one agent route. The fake goal connects all exits to a single
sentinel, making every reachable cell point toward a suitable exit. Directions
include eight planar directions and portal transitions \texttt{DIR\_IN} and
\texttt{DIR\_OUT}.

\begin{table}[H]
\centering
\begin{tabular}{|p{0.31\textwidth}|p{0.61\textwidth}|}
\hline
\textbf{State} & \textbf{Role} \\
\hline
Fake goal &
Sentinel outside the real grid used as a common target for all emergency exits. \\
\hline
Direction map &
Per-floor arrays of \texttt{domain.Direction} values sent through
\texttt{FlowSub}. \\
\hline
World subscription &
Tracks changed cells and queues signaling nodes whose local flow may need
refreshing. \\
\hline
\texttt{flowPublisher} &
Single publisher for emergency flow snapshots. A new subscription replaces the
previous publisher. \\
\hline
\end{tabular}
\caption{Main state of the emergency flow field.}
\end{table}

\subsection{Environment Updates and Signaling}

Before computing or updating the flow field, \texttt{applyEnvironmentObjects}
queries all oriented objects and refreshes object layers for all floors. Each
floor uses the region linked to that layer, clears previous object costs,
rasterizes relevant objects, and calls \texttt{UpdateFinal}. Changed cells are
registered by the world's subscription.

\begin{table}[H]
\centering
\begin{tabular}{|p{0.32\textwidth}|p{0.60\textwidth}|}
\hline
\textbf{Parameter or dependency} & \textbf{Emergency effect} \\
\hline
\texttt{LightingHystheresis} &
Minimum cost-change threshold before a signaling node direction is updated. \\
\hline
\texttt{PriorityPathCellsWidth} &
Width, in grid cells, used to rasterize preferred routes onto the grid. \\
\hline
\texttt{AllObjectsOriented} &
Source of the objects used to refresh dynamic costs on all floors. \\
\hline
\texttt{SignalingSetDirection} &
External operation that writes the selected direction into a signaling node. \\
\hline
\texttt{updateTicksPerSecond} &
Frequency at which the emergency goroutine refreshes the environment, repairs
planning state, and publishes changes when needed. \\
\hline
\end{tabular}
\caption{Configuration and dependencies relevant to emergency routing.}
\end{table}

\subsection{Preferred Routes and Restoration}

Preferred routes arrive as per-floor graphs. \texttt{ApplyPreferencePaths}
searches the waypoints of each edge and rasterizes a line on the corresponding
floor grid. This line applies a negative mask to the base layer so the flow field
favors those corridors. Malformed edges with missing waypoints are ignored. If
an edge omits its weight, the implementation uses the default value \texttt{1}.

When the emergency stops, \texttt{RemovePreferencePaths} calls
\texttt{ResetBase} on every floor to restore normal traversable-space cost while
keeping object layers current. Restoration must happen both on ordinary stop and
during service shutdown.

\begin{table}[H]
\centering
\begin{tabular}{|p{0.32\textwidth}|p{0.60\textwidth}|}
\hline
\textbf{Responsibility} & \textbf{Code path} \\
\hline
Apply preferred routes &
The start handler calls \texttt{ApplyPreferencePaths} before creating the flow
field. \\
\hline
Remove routes after startup error &
If goal conversion or initialization fails, the handler removes preference masks
before returning the error. \\
\hline
Remove routes on stop &
\texttt{completeEmergencyStopLocked} calls \texttt{RemovePreferencePaths} and
marks the emergency stopped. \\
\hline
Remove routes during shutdown &
\texttt{stopCurrentEmergency} joins the emergency goroutine and performs the same
cleanup even if shutdown races with a client stop request. \\
\hline
\end{tabular}
\caption{Emergency cleanup responsibilities.}
\end{table}

\subsection{Validation of Emergency Changes}

Emergency changes should cover startup, blocked goals, simultaneous start
rejection, flow publication, object updates, signaling, preferred routes, and
ordered stop. If a change affects signaling, include tests for hysteresis,
direction selection, and missing or invalid nodes.

\section{Public API, NATS Subjects, and Dispatcher}

The public API has two equivalent views: typed Go wrappers and raw NATS
messages. The wrappers live in \texttt{pkg/pathfinding} and
\texttt{pkg/emergency}; the raw subject names live in \texttt{pkg/subjects}. The
dispatcher owns the NATS subscriptions and translates messages into internal
grid, world, communicator, and flow-field operations.

\subsection{Boundary Between Public Clients and Internal Handlers}

\begin{table}[H]
\centering
\begin{tabular}{|p{0.29\textwidth}|p{0.29\textwidth}|p{0.32\textwidth}|}
\hline
\textbf{Public operation} & \textbf{Internal payload or handler} &
\textbf{Internal effect} \\
\hline
\texttt{Pathfinding.} \newline \texttt{AgentFindPath} &
\texttt{AgentFindPathParams} &
Converts goal, creates an agent communicator, and registers it in the dispatcher. \\
\hline
\texttt{Pathfinding.} \newline \texttt{AgentNaivePathCost} &
\texttt{AgentNaivePath-} \newline \texttt{CostParams} &
Computes cost without creating movement state. \\
\hline
\texttt{Pathfinding.} \newline \texttt{AgentPathSub} &
\texttt{IdentifierParams} and \texttt{snapshotstream.} \newline
\texttt{Envelope} &
Streams route cell states through the request reply subject. \\
\hline
\texttt{Emergency.Start} &
\texttt{EmergencyPathParams} &
Applies preferred routes, validates goals, starts the flow field, and updates
emergency lifecycle state. \\
\hline
\texttt{Emergency.FlowSub} &
\texttt{snapshotstream.} \newline \texttt{Envelope} &
Streams emergency direction maps per floor. \\
\hline
\end{tabular}
\caption{Public operations and internal dispatcher paths.}
\end{table}

\subsection{Subject Registration and Concurrency}

\texttt{app.Run} subscribes all public subjects. Most handlers execute directly
on NATS callbacks. \texttt{AgentFindPath}, \texttt{BlockingWait}, and
\texttt{MoveFMeters} are wrapped by bounded worker handlers because they can
block long enough to occupy the callback path.

\begin{table}[H]
\centering
\begin{tabular}{|p{0.29\textwidth}|p{0.23\textwidth}|p{0.38\textwidth}|}
\hline
\textbf{Subject group} & \textbf{Worker limit} & \textbf{Reason} \\
\hline
\texttt{pathfinding.} \newline \texttt{agent\_find\_path} &
\texttt{findpath-} \newline \texttt{handlerworkers} &
Route initialization can block while preparing the pathfinding run. \\
\hline
\texttt{agent\_comm.} \newline \texttt{blocking\_wait} &
\texttt{blockingwait-} \newline \texttt{handlerworkers} &
The request waits until an agent completes. \\
\hline
\texttt{agent\_comm.} \newline \texttt{move\_f\_meters} &
\texttt{movefmeters-} \newline \texttt{handlerworkers} &
The request waits until a distance is consumed or the route ends. \\
\hline
Other request subjects &
No extra worker wrapper &
They are expected to finish quickly or stream through a separate goroutine. \\
\hline
\end{tabular}
\caption{NATS handler concurrency model.}
\end{table}

\subsection{Agent Session Lifecycle}

A new \texttt{AgentFindPath} call for an agent replaces the previous run for
that same identifier. The old communicator is terminated. Completed results are
retained for a short interval so \texttt{IsMoving}, \texttt{BlockingWait}, and
\texttt{ExitError} can still respond after the active communicator has left the
dispatcher table. During shutdown, \texttt{stopActiveAgents} prevents new
registrations, terminates all active communicators, waits for them, and clears
the tables.

\subsection{Subscription Flows}

Path and flow subscriptions use request/reply setup but stream multiple messages
to the reply subject. The handler returns immediately and a goroutine owns the
reply stream. Messages contain \texttt{done=false} while data continues and a
final \texttt{done=true} message before the channel closes.

Small snapshots keep the legacy envelope:

\begin{lstlisting}
{"done": false, "cells": {"0": [0, 1, 4]}}
\end{lstlisting}
Large snapshots use the shared \texttt{pkg/snapshotstream} codec. The service
serializes the per-floor map, compresses it with \texttt{gzip}, and publishes a
sequence of chunk envelopes:

\begin{lstlisting}
{
  "encoding": "gzip+json",
  "snapshotId": "42",
  "chunkIndex": 0,
  "chunkCount": 3,
  "data": "..."
}
\end{lstlisting}

\texttt{chunkIndex} is zero-based within one \texttt{snapshotId}. Because
\texttt{data} is a Go \texttt{[]byte} encoded as JSON, non-Go clients must first
base64-decode that field for each chunk. Consumers then join all chunks for the
same snapshot, decompress the combined payload, and decode the JSON as the same
per-floor map otherwise carried by \texttt{cells}. The Go wrappers use
\texttt{snapshotstream.Assembler} and expose only reconstructed snapshots on
\texttt{AgentPathSub} and \texttt{FlowSub}. Service handlers should use
\texttt{snapshotstream.MarshalSnapshot} for data messages and
\texttt{snapshotstream.MarshalDone} for terminal messages.

Replacing an agent path publisher closes the previous publisher for that
communicator. Replacing the emergency flow publisher closes the previous global
flow publisher. Clients must be ready for stream closure and should resubscribe
if they need to become the active observer again.

\subsection{Public Errors}

Domain errors are registered with the shared message package so raw NATS clients
can receive stable error codes. Go wrappers additionally wrap failures with
operation and NATS-phase sentinels such as \texttt{ErrNATSRequest},
\texttt{ErrNATSResponse}, \texttt{ErrNATSSubscription}, and
\texttt{ErrNATSPublish}.

\begin{table}[H]
\centering
\begin{tabular}{|p{0.36\textwidth}|p{0.55\textwidth}|}
\hline
\textbf{Public error} & \textbf{Typical situation} \\
\hline
\texttt{ErrAgentCommNotFound} &
A communicator operation references an identifier without an active or retained
session. \\
\hline
\texttt{ErrAgentNoPath} &
No traversable route exists from the current agent position to the goal. \\
\hline
\texttt{ErrAgentTerminated} &
The route was explicitly terminated. \\
\hline
\texttt{ErrAgentStillRunning} &
\texttt{ExitError} was queried while the route is still active. \\
\hline
\texttt{ErrAgentExitedWithMetersLeft} &
A meter movement request could not consume the full requested distance because
the route ended first. \\
\hline
\texttt{ErrEmergencyStillRunning} &
An emergency start was requested while another emergency was active or stopping. \\
\hline
\texttt{ErrDispatcherShuttingDown} &
The dispatcher is draining and no longer accepts new agent registrations. \\
\hline
\end{tabular}
\caption{Domain errors that are part of the public API.}
\end{table}

\subsection{Validation of Public API Changes}

When adding or changing public API behavior, update all related contracts
together: subject constants, payload type, dispatcher handler, Go wrapper,
integration tests, and user manual. Include malformed-payload or missing-session
cases for handlers that accept external input.

\section{Configuration and Embedded Execution}

Configuration affects both the standalone executable and embedded execution. It
defines NATS connectivity, logging, grid shape, layers, portals, handler worker
limits, geotruth pool sizes, and algorithm parameters.

\subsection{Configuration Sources}

\texttt{main.go} calls \texttt{LoadConfig(nil)}. This creates a Viper instance,
registers defaults from \texttt{DefaultConfig}, loads
\texttt{config/config.yaml}, and allows \texttt{PTHFSERVICE\_*} environment
variables to override keys. Tests, tools, and benchmarks can pass an explicit
YAML path.

\begin{table}[H]
\centering
\begin{tabular}{|p{0.28\textwidth}|p{0.22\textwidth}|p{0.39\textwidth}|}
\hline
\textbf{Source} & \textbf{Precedence} & \textbf{Use} \\
\hline
\texttt{DefaultConfig} &
Lowest &
Provides local defaults for NATS, grid, maps, algorithm parameters, and worker
limits. \\
\hline
\texttt{config/config.yaml} or explicit path &
Middle &
Project or scenario configuration. \\
\hline
\texttt{PTHFSERVICE\_*} &
Highest &
Operational override. Dotted keys become underscores, for example
\texttt{PTHFSERVICE\_APP\_SERVERURL}. \\
\hline
\end{tabular}
\caption{Configuration sources and precedence.}
\end{table}

\subsection{Public Configuration Structure}

\begin{table}[H]
\centering
\begin{tabular}{|p{0.28\textwidth}|p{0.63\textwidth}|}
\hline
\textbf{Type} & \textbf{Responsibility} \\
\hline
\texttt{AppConfig} &
NATS URL and logging destinations. \\
\hline
\texttt{GridConfig} &
Layers, portals, rows, columns, pixels per cell, meters per cell, and wall aura. \\
\hline
\texttt{PathfindingConfig} &
Shared free-movement height plus the agent, emergency, and geotruth blocks. \\
\hline
\texttt{AgentConfig} &
Vision radius, real-position update cadence, reusable-world limit, sparse value
storage, NATS worker limits, and debug output. \\
\hline
\texttt{EmergencyConfig} &
Signaling hysteresis, preferred-route width in cells, and debug output. \\
\hline
\texttt{GeotruthConfig} &
Internal NATS connection pool sizes for geotruth queries and publishes. \\
\hline
\end{tabular}
\caption{Public configuration groups.}
\end{table}

\subsection{Dependency Injection}

\begin{table}[H]
\centering
\begin{tabular}{|p{0.23\textwidth}|p{0.31\textwidth}|p{0.36\textwidth}|}
\hline
\textbf{Field} & \textbf{Interface} & \textbf{Purpose} \\
\hline
\texttt{Ex} &
\texttt{ExternalSystem} &
Object traversal costs, floor heights, and emergency signaling operations. \\
\hline
\texttt{Qu} &
\texttt{GeoQuery} &
Object data, local vision, all oriented objects, and floor/region lookup. \\
\hline
\texttt{Pu} &
\texttt{GeoPublish} &
Movement publication and commit acknowledgement for real object positions. \\
\hline
\end{tabular}
\caption{Injected dependencies and subsystem use.}
\end{table}

\texttt{DefaultDependencies} is useful for local runs and tests because it
creates geotruth NATS pools and a simulated \texttt{ExternalSystem}. In a real
installation, the integrator should provide implementations that match the
deployment's geotruth and signaling services.

\subsection{Standalone and Embedded Execution}

The standalone executable and embedded API converge in \texttt{app.Run}. In
standalone mode, \texttt{main.go} loads the file, creates default dependencies,
and waits for OS signals. In embedded mode, the caller constructs
\texttt{Config}, supplies dependencies, calls \texttt{embedded.Run}, and later
calls \texttt{Shutdown} on the returned dispatcher.

\begin{table}[H]
\centering
\begin{tabular}{|p{0.25\textwidth}|p{0.30\textwidth}|p{0.35\textwidth}|}
\hline
\textbf{Responsibility} & \textbf{Standalone} & \textbf{Embedded} \\
\hline
Configuration &
\texttt{LoadConfig(nil)} from \texttt{config/config.yaml}. &
Caller can use defaults, explicit load, or construct a configuration. \\
\hline
Dependencies &
\texttt{embedded.} \newline \texttt{DefaultDependencies}. &
Caller should usually provide production dependencies. \\
\hline
Lifecycle &
OS signal triggers shutdown. &
Caller owns cancellation and must call \texttt{Shutdown}. \\
\hline
\end{tabular}
\caption{Standalone and embedded execution responsibilities.}
\end{table}

\subsection{Validation of Configuration Changes}

For a new configuration key, add the public field, \texttt{mapstructure} tag,
default value, documentation, and tests for defaults, YAML loading, and
environment override. If the value affects runtime behavior, add a functional
test in the subsystem that consumes it.

\section{External Integrations and Test Doubles}

The pathfinding service depends on external systems but keeps them behind small
interfaces. Internal packages should not create real external clients by
themselves. They receive dependencies prepared by the standalone startup,
embedded API, or tests.

\subsection{Dependency Contracts}

\begin{table}[H]
\centering
\begin{tabular}{|p{0.22\textwidth}|p{0.28\textwidth}|p{0.40\textwidth}|}
\hline
\textbf{Contract} & \textbf{Primary users} & \textbf{Notes} \\
\hline
\texttt{ExternalSystem} &
Agent movement, emergency signaling, object costs. &
Must propagate errors with enough context to diagnose external failures. \\
\hline
\texttt{GeoQuery} &
Coordinate conversion, agent local vision, emergency object refresh. &
Region names returned by \texttt{RegionFromPoint} must match configured layer
names. \\
\hline
\texttt{GeoPublish} &
Agent real-position publishing. &
Returns a commit acknowledgement. Resource ownership belongs to the concrete
implementation or the pool that wraps it. \\
\hline
\end{tabular}
\caption{External dependency contracts.}
\end{table}

\subsection{Use by Subsystem}

Agent planning uses \texttt{GeoQuery.ObjectData}, \texttt{RegionFromPoint},
\texttt{NearbyObjectsOf}, \newline \texttt{ObjectTraversalCost},
\texttt{HeightAtFloorPoint}, and \texttt{GeoPublish.UpdateObjectPosition}.
Emergency routing uses \texttt{AllObjectsOriented}, \texttt{LightNodeRegex}, and
\newline \texttt{SignalingSetDirection}. Startup resolves the available signaling nodes
once when the dispatcher is created.

\subsection{Embedded Execution and Real Services}

Embedded execution can be used with the default dependencies or with real
services. The default dependency path still connects to NATS and creates
geotruth query/publish pools, but it uses a simulated external system for domain
costs and signaling. For production, provide dependencies that implement the
public interfaces and close any external resources they own.

\subsection{Test Doubles}

\texttt{internal/mock} provides test doubles for the domain and geotruth parts
needed by the service. They are used for unit, component, integration, and
performance tests where deterministic behavior is more important than full
external-system fidelity.

\begin{table}[H]
\centering
\begin{tabular}{|p{0.25\textwidth}|p{0.36\textwidth}|p{0.29\textwidth}|}
\hline
\textbf{Double} & \textbf{Simulated behavior} & \textbf{Main use} \\
\hline
\texttt{mock.GeoTruth} &
In-memory object store, lookup by ID, local object queries, and oriented object
sets. &
Agent and emergency tests. \\
\hline
\texttt{mock.ExternalSystem} &
Object traversal costs, floor heights, and signaling methods. &
Algorithm and integration tests. \\
\hline
\texttt{internal/} \newline \texttt{pathfinding/testing} scenarios &
Reusable map and route cases. &
Black-box pathfinding behavior. \\
\hline
\end{tabular}
\caption{Main test doubles and scenarios.}
\end{table}

\section{System Extension and Technical Limits}

The system supports local changes, but some resources are process-wide:
configuration, dependencies, and the global world. An extension must not
initialize unrelated worlds or modify global configuration while agents or
emergencies are active unless it also defines and tests a new synchronization
strategy.

\subsection{Extension Points}

\begin{table}[H]
\centering
\begin{tabular}{|p{0.23\textwidth}|p{0.34\textwidth}|p{0.32\textwidth}|}
\hline
\textbf{Extension} & \textbf{Contracts to update} & \textbf{Minimum validation} \\
\hline
New NATS subject or public operation &
\texttt{pkg/subjects}, dispatcher payload, handler registration, Go wrapper,
and user manual. &
Integration test, invalid-payload case, and public-error review. \\
\hline
Snapshot stream envelope &
\texttt{pkg/snapshotstream}, dispatcher publication helpers, Go stream
assemblers, raw-NATS documentation, and user manual. &
Codec tests plus public \texttt{AgentPathSub} and \texttt{FlowSub} round trips
for chunked snapshots. \\
\hline
New agent communicator command &
Public method, subject, handler, internal communicator method, and behavior for
finished or missing agents. &
Tests with active agent, terminated agent, cancellation, and post-completion
calls. \\
\hline
New configuration parameter &
Public field, \texttt{mapstructure} tag, default, file/env loading, and
documentation. &
Configuration-load test and subsystem functional test. \\
\hline
New external dependency &
Public interface method, real implementation, test double, and subsystem call. &
Double-based test, error propagation test, and integration case when NATS or
geotruth is involved. \\
\hline
New algorithm or search variant &
Internal input boundary, state types, global-grid interaction, and result
publication. &
Algorithm unit tests, no-path cases, multi-floor cases if applicable, and public
API test when externally observable. \\
\hline
\end{tabular}
\caption{Main extension points and validation requirements.}
\end{table}

\subsection{Current Technical Limits}

\begin{table}[H]
\centering
\begin{tabular}{|p{0.27\textwidth}|p{0.31\textwidth}|p{0.31\textwidth}|}
\hline
\textbf{Limit} & \textbf{Current behavior} & \textbf{Implication} \\
\hline
Process-wide world &
The global world and layer translation are shared. &
Multiple independent worlds in one process require redesign. \\
\hline
Static signaling set &
Light nodes are resolved when the dispatcher starts. &
Nodes appearing or disappearing during an emergency are future work. \\
\hline
Grid-based movement &
Public coordinates are meters, but algorithms operate on rasterized cells. &
Precision and policy rules depend on grid resolution. \\
\hline
Preferred-route restoration &
Emergency preferred routes modify global base costs and are reset on stop. &
Unexpected shutdown paths must preserve cleanup. \\
\hline
Per-agent local world &
Agents copy or link local dynamic state per route. &
Large maps and many agents may require paging or higher-level guidance. \\
\hline
\end{tabular}
\caption{Current technical limits and their design intent.}
\end{table}

\subsection{Checklist Before Accepting an Extension}

Before accepting a shared change, identify whether it affects configuration,
grid/rasterization, algorithms, public API, external integrations, or lifecycle.
Then confirm documentation, tests, and public wrappers are updated together. If
a limitation remains, document it here or in the project conclusions rather than
letting it remain implicit.

\section{Tests, Benchmarks, and Technical Validation}

The repository uses Go's \texttt{testing} package. Tests live beside the
packages they exercise, with additional integration-style tests under
\texttt{internal/integration} and shared behavior scenarios under
\texttt{internal/pathfinding/testing}. Heavy diagnostics and benchmarks are kept
separate from ordinary tests.

\subsection{Test Map by Responsibility}

\begin{table}[H]
\centering
\begin{tabular}{|p{0.24\textwidth}|p{0.31\textwidth}|p{0.33\textwidth}|}
\hline
\textbf{Area} & \textbf{Typical packages} & \textbf{Purpose} \\
\hline
Grid and raster &
\texttt{internal/grid}, \texttt{internal/raster}, \texttt{pkg/tuning} &
Coordinate conversion, cost layers, image rasterization, wall aura, portals,
subscriptions, and visualization helpers. \\
\hline
Agent algorithm &
\texttt{internal/} \newline \texttt{pathfinding/agent} &
D* Lite behavior, local changes, movement, communicator commands, and
multi-floor paths. \\
\hline
Emergency algorithm &
\texttt{internal/} \newline \texttt{pathfinding/emergency} &
Flow-field planning, preferred routes, signaling, environment updates, and stop
cleanup. \\
\hline
Public API &
\texttt{pkg/pathfinding}, \texttt{pkg/emergency}, \texttt{internal/app},
\texttt{internal/integration} &
NATS request/response behavior, wrappers, errors, lifecycle, and shutdown. \\
\hline
Configuration and embedded &
\texttt{internal/config}, \texttt{embedded}, integration tests &
Default values, YAML loading, environment overrides, embedded startup, and
dependency cleanup. \\
\hline
\end{tabular}
\caption{Test responsibilities by subsystem.}
\end{table}

\subsection{Unit and Component Tests}

Unit tests validate small, deterministic behavior: priority queues, raster
primitives, grid cost math, coordinate conversion, portal lookup, paged value
storage, and public wrapper error classification. Component tests exercise
larger behavior such as D* Lite replanning, agent communicators, flow-field
updates, and dispatcher shutdown.

\subsection{API Boundary Error Tests}

The API tests include representative malformed JSON, missing communicator,
dispatcher-shutdown, emergency-already-running, idempotent stop, and post-agent
completion cases. The goal is scenario-driven validation, not duplicating every
malformed payload variant for every subject.

\subsection{Public API Integration Tests}

Integration tests start NATS and the service path, then use public clients or raw
subjects to validate observable behavior. They check direct routes, multi-floor
routes, subscriptions, movement commands, emergency start/stop, and error
propagation through the same public surface that external consumers use.

\subsection{Benchmarks and Heavy Tests}

Benchmarks are used for diagnosis and scalability evidence. Ordinary benchmarks
measure local structures or algorithms. Heavy benchmarks require explicit
environment variables and should be run only when evaluating capacity or
full-stack behavior.

\begin{table}[H]
\centering
\begin{tabular}{|p{0.27\textwidth}|p{0.30\textwidth}|p{0.31\textwidth}|}
\hline
\textbf{Benchmark group} & \textbf{Focus} & \textbf{Use} \\
\hline
Grid/raster benchmarks &
Map loading, portal lookup, world creation. &
After changing rasterization or world data structures. \\
\hline
Agent benchmarks &
D* Lite planning, local cost updates, movement capacity. &
After changing agent planning, movement cadence, or virtual worlds. \\
\hline
Emergency benchmarks &
Flow-field planning, object refresh, preferred-route application. &
After changing emergency recomputation or signaling behavior. \\
\hline
API capacity benchmarks &
NATS wrappers, handler workers, geotruth pool settings. &
When evaluating full-stack concurrency. \\
\hline
\end{tabular}
\caption{Benchmark groups and when to use them.}
\end{table}

Typical heavy benchmark commands include:

\begin{Verbatim}[breaklines=true,fontsize=\small]
$env:PATHFINDING_BENCH_HEAVY='1'; go test ./internal/performance -run ^$ -bench BenchmarkSynthetic_AgentCapacity -benchtime=1x -count=7 -benchmem
$env:PATHFINDING_BENCH_HEAVY='1'; $env:PATHFINDING_BENCH_REAL='1'; go test ./internal/performance -run ^$ -bench BenchmarkReal_AgentCapacity -benchtime=1x -count=7 -benchmem
$env:PATHFINDING_BENCH_HEAVY='1'; go test ./internal/pathfinding/emergency -run ^$ -bench BenchmarkFlowfield_MovingObjectsCapacity -benchtime=1x -count=7 -benchmem
\end{Verbatim}

When using \texttt{-count=7}, summarize latency with the median and treat a point
as stable only when all seven runs satisfy the benchmark criterion. Partial
success should be described as unstable, not as guaranteed capacity.

Agent capacity benchmarks distinguish \texttt{RawStep}, \texttt{ExternalTick},
and \texttt{RealTime}. \newline \texttt{ExternalTick} best represents simulations whose
clock is controlled by another system: agents start stopped, external ticks ask
each active agent to move a meter distance, and leftover distance is accumulated.
\texttt{RealTime} measures autonomous movement where the service controls
cadence.

Emergency moving-object capacity measures a synchronous emergency tick: read
oriented objects, rasterize them, apply changed cells, and repair the flow field.
The main axis is \texttt{moving\_objects}; \texttt{p95\_tick\_ms} is compared
with the target tick budget, while \texttt{missed\_ticks} remains diagnostic.

\subsection{Coverage, Execution, and Change Validation}

The main validation command is:

\begin{lstlisting}
make test
\end{lstlisting}

This target runs \texttt{go test} with verbose output, local package coverage,
and \texttt{coverage.out}. Local coverage shows what each package's own tests
exercise, but it does not attribute code executed by tests from other packages.
For complete cross-package coverage, use:

\begin{lstlisting}
make test-cover-full
\end{lstlisting}

This target runs all module tests with \texttt{-coverpkg=./...} and produces
\texttt{coverage-full.out}. Standard Go coverage tools can open both profiles.

\begin{table}[H]
\centering
\begin{tabular}{|p{0.28\textwidth}|p{0.28\textwidth}|p{0.31\textwidth}|}
\hline
\textbf{Validation type} & \textbf{Command} & \textbf{When to use it} \\
\hline
Full tests &
\texttt{make test} &
Before closing changes to algorithms, grid, API, configuration, or lifecycle. \\
\hline
Full coverage &
\texttt{make test-cover-full} &
When measuring all code exercised by any test in the module. \\
\hline
Specific package &
\texttt{go test ./internal/...} or a specific package &
During local iteration on a small area. \\
\hline
Local coverage HTML &
\texttt{make cover-html} &
After \texttt{make test}. \\
\hline
Cross-package coverage HTML &
\texttt{make cover-full-html} &
After \texttt{make test-cover-full}. \\
\hline
Benchmarks &
\texttt{go test -bench=.} on the affected package &
When changing data structures, rasterization, planning, or resource reuse. \\
\hline
\end{tabular}
\caption{Recommended technical validation commands.}
\end{table}

\section{Maintenance, Diagnosis, and Operational Problems}

Maintenance requires observing logs, external resources, and active route
lifecycle together. Most visible API errors come from inconsistent
configuration, unavailable NATS or geotruth, coordinates that do not map to the
grid, missing layers, disconnected routes, or incomplete cleanup during stop.

\subsection{Logs and Diagnosis}

Internal logging uses a \texttt{[PATHFINDING]} prefix, date, time, and source
file. Configuration chooses standard error, a file, both, or neither. Disabling
all output can be useful for tests but makes production diagnosis difficult.

\begin{table}[H]
\centering
\begin{tabular}{|p{0.25\textwidth}|p{0.32\textwidth}|p{0.32\textwidth}|}
\hline
\textbf{Area} & \textbf{Observable signal} & \textbf{First check} \\
\hline
Startup &
Panic or error before subjects are registered. &
Configuration, map paths, NATS URL, and default dependencies. \\
\hline
Rasterization &
Complex-image warning or failure opening a map path. &
Image path relative to the execution directory and configured dimensions. \\
\hline
Portals &
Portal dropped because of invalid cost, endpoint, or same layer. &
Cost, source/destination layers, and meter coordinates. \\
\hline
NATS API &
Client timeout or no response. &
Same NATS server, registered subject, and service startup logs. \\
\hline
Agents &
Communicator error, no path, or real-position update failure. &
Agent existence in geotruth, valid goal floor, connectivity, and publisher
acknowledgement. \\
\hline
Emergency &
Start failure, unchanged flow, or no signaling update. &
Goals, preferred routes, light nodes, and oriented-object data. \\
\hline
\end{tabular}
\caption{Initial diagnosis points by symptom.}
\end{table}

\subsection{Graceful Shutdown}

Normal shutdown enters through \texttt{Shutdown}. In the standalone executable,
this follows \texttt{SIGINT} or \texttt{SIGTERM}. In embedded mode, the caller
must invoke it. The order matters because agent goroutines, emergency
goroutines, publishers, and NATS subscriptions may be active simultaneously.

\begin{table}[H]
\centering
\begin{tabular}{|p{0.27\textwidth}|p{0.29\textwidth}|p{0.31\textwidth}|}
\hline
\textbf{Phase} & \textbf{Resource} & \textbf{Risk if it fails} \\
\hline
Stop agents &
Active communicators and retained results. &
Blocked goroutines, movements after shutdown, or lost terminal results. \\
\hline
Stop emergency &
Quit channel, done channel, preferred routes. &
Base costs not restored or flow publisher left open. \\
\hline
Drain NATS &
Dispatcher subscriptions and pending replies. &
Clients miss final replies or callbacks run while state is being released. \\
\hline
Cancel context &
Root service context. &
Internal operations wait on already closed resources. \\
\hline
Close external resources &
Connections, clients, or descriptors owned by dependencies. &
Resource leaks in the embedding process. \\
\hline
\end{tabular}
\caption{Relevant phases of graceful shutdown.}
\end{table}

\subsection{Common Problems}

\begin{table}[H]
\centering
\begin{tabular}{|p{0.27\textwidth}|p{0.31\textwidth}|p{0.30\textwidth}|}
\hline
\textbf{Problem} & \textbf{Likely cause} & \textbf{Initial action} \\
\hline
Timeout on a public call &
NATS unreachable, subject not registered, or service not started. &
Check URL, service startup, and subject registration logs. \\
\hline
Agent not found &
Identifier missing in geotruth or completed result expired. &
Verify object registration and distinguish agent ID from active session state. \\
\hline
No route &
Goal outside grid, missing layer, walls without portals, or unreachable cost. &
Check converted coordinates, layers, portals, and rasterized map. \\
\hline
Movement without real update &
Publishing dependency failure or unsuitable update cadence. &
Check \texttt{GeoPublish}, commit acknowledgement, floor height, and update
configuration. \\
\hline
Emergency does not start &
Another emergency is active, goal is invalid, or flow creation fails. &
Check dispatcher emergency state, goals, and preferred-route restoration on
failure. \\
\hline
Signaling does not change &
Light nodes not resolved, hysteresis too high, or dominant direction unchanged. &
Check node query, computed directions, and hysteresis. \\
\hline
\end{tabular}
\caption{Common operational problems and first technical action.}
\end{table}

\subsection{Maintenance Checklist}

Before changing a shared area, identify whether the change affects
configuration, grid, algorithm, public API, or lifecycle. Then review external
contracts and run the minimum validation for the affected subsystem.

\begin{table}[H]
\centering
\begin{tabular}{|p{0.34\textwidth}|p{0.58\textwidth}|}
\hline
\textbf{Review} & \textbf{Criterion} \\
\hline
Effective configuration &
The changed key has a default, file loading behavior, and documented effect. \\
\hline
External resources &
Queries, publications, and external-system calls have timeouts and contextual
errors. \\
\hline
Resource ownership &
The owner of subscriptions, channels, worlds, publishers, and connections is
clear. \\
\hline
Public compatibility &
If an error, subject, JSON field, or public type changes, update the user manual
and integration tests. \\
\hline
Diagnostic logs &
Logs help locate failure phase without leaking sensitive data or flooding normal
output. \\
\hline
Tests and benchmarks &
Run subsystem validation and compare benchmarks before and after performance
changes. \\
\hline
\end{tabular}
\caption{Technical checklist before closing maintenance changes.}
\end{table}

\subsection{Validation After Incidents}

A fixed incident should leave a reproducible test or check. Configuration
failures need loading or invalid-value tests. Concurrency failures need tests for
the observed race. Integration failures should exercise the public path that
failed.

\begin{table}[H]
\centering
\begin{tabular}{|p{0.34\textwidth}|p{0.58\textwidth}|}
\hline
\textbf{Incident} & \textbf{Recommended validation} \\
\hline
Startup error &
Configuration or embedded-start test, plus logging check if the failure phase
was unclear. \\
\hline
NATS timeout or error &
Application or integration test confirming subject registration, response, and
shutdown without pending goroutines. \\
\hline
Incorrect route &
Algorithm unit test and integration test if the difference was visible through
the public API. \\
\hline
Contaminated global state &
Reuse or cleanup test for virtual worlds, local portals, preferred routes, or
completed results. \\
\hline
Shutdown error &
\texttt{Shutdown} test with active agents, active emergency, or simultaneous
client stop. \\
\hline
Performance regression &
Benchmark of the affected component with the same map and configuration before
and after the correction. \\
\hline
\end{tabular}
\caption{Recommended validation after operational incidents.}
\end{table}

\begin{thebibliography}{9}
\addcontentsline{toc}{chapter}{Bibliography}

\bibitem{koenig-likhachev-dstar-lite}
S. Koenig and M. Likhachev, \enquote{D* Lite}, in \emph{Proceedings of the
Eighteenth National Conference on Artificial Intelligence (AAAI-02)}, Edmonton,
Alberta, Canada, 2002, pp. 476--483. Available:
\url{https://cdn.aaai.org/AAAI/2002/AAAI02-072.pdf}. Accessed: June 15, 2026.

\bibitem{foley1995computergraphics}
J. D. Foley, A. van Dam, S. K. Feiner, and J. F. Hughes, \emph{Computer
Graphics: Principles and Practice in C}, 2nd ed. Reading, MA, USA:
Addison-Wesley, 1995.

\bibitem{pineda1988polygon}
J. Pineda, \enquote{A Parallel Algorithm for Polygon Rasterization}, \emph{ACM
SIGGRAPH Computer Graphics}, vol. 22, no. 4, pp. 17--20, 1988, doi:
10.1145/54852.378457.

\end{thebibliography}

\end{document}
